% ============================================================
% Chapter 1 -- Introduction
% ============================================================
\chapter{Introduction}
\label{ch:introduction}

\section{Background}
\label{sec:background}

Dhaka is one of the most crowded cities in the world. Every day, millions of
people travel to offices, universities and markets using buses, CNG
auto-rickshaws, rickshaws and ride-hailing services such as Uber and Pathao.
Public buses are cheap but crowded, uncomfortable and often unsafe, especially
for women. Ride-hailing services are comfortable, but a private car or CNG ride
is too expensive for most students and working people to use every day.

At the same time, most private cars and CNGs on the road carry only one
passenger, even though they have space for two or three. Many of these
passengers are travelling from the same area to the same destination at the
same time. If these people could share one vehicle and split the fare, each of
them would pay much less, and fewer vehicles would be needed on the road.

Ride pooling (also called ride sharing or carpooling) is the idea of matching
passengers who are going the same way so that they can travel together in one
vehicle. This project, \textbf{RidePool}, is a ride-pooling platform built for
Dhaka. It consists of a mobile app for riders (CarPoolApp), a mobile app for
drivers (DriverApp) and a backend server that matches riders into shared rides,
calculates fair fares and manages payments.

\section{Motivation}
\label{sec:motivation}

The main motivation behind RidePool is \textbf{cost saving}. Existing platforms
in Bangladesh are designed mainly for single-passenger trips, and their pooling
options are either limited or not available. We wanted to build a system where
sharing a ride is the main purpose, not an extra feature.

Other reasons that motivated this project are:

\begin{itemize}
  \item \textbf{High transport cost:} Daily ride-hailing is not affordable for
        students and low- to middle-income workers. Splitting a fare between
        two or three people makes a comfortable ride affordable.
  \item \textbf{Traffic congestion:} Fewer vehicles carrying more passengers
        can help reduce traffic and fuel usage in the city.
  \item \textbf{Safety and comfort:} Many passengers, especially women, feel
        uncomfortable sharing a ride with strangers. A pooling system with a
        female-only option and a trusted-companion feature can make shared
        rides feel safer.
  \item \textbf{Local needs:} Most ride-sharing solutions are designed for
        foreign markets. Dhaka needs a solution that supports local vehicles
        like CNGs, local currency (BDT) and local payment habits such as mobile
        banking and cash.
\end{itemize}

\section{Problem Discussion}
\label{sec:problem}

Although ride pooling sounds simple, building it in practice raises several
problems:

\begin{enumerate}
  \item \textbf{Finding the right co-passengers:} The system must quickly find
        riders whose pickup points are close to each other and whose
        destinations lie in the same direction. Comparing every rider with
        every other rider is slow and does not scale.
  \item \textbf{Fair fare sharing:} When several people share a ride, each
        person should pay less than a solo ride, and the fare should be
        calculated in a clear and fair way.
  \item \textbf{Waiting time:} A rider does not want to wait a long time for
        other passengers to join. The system must form a pool quickly or tell
        the rider when no pool is available.
  \item \textbf{Safety and trust:} Passengers are sharing a vehicle with
        strangers. The system needs gender-based preferences and a way to ride
        with trusted people.
  \item \textbf{Cancellations:} In a shared ride, one person cancelling affects
        everyone else in the pool. Frequent cancellations must be discouraged.
  \item \textbf{Concurrency:} Many riders may try to join the same pool at the
        same moment, and many drivers may try to accept the same pool. The
        system must never overfill a vehicle or assign two drivers to one pool.
  \item \textbf{Planning ahead:} Many daily trips (office, university) are
        known in advance, but most services only support instant booking.
\end{enumerate}

RidePool is designed to solve these problems in one integrated system.

\section{Project Aims}
\label{sec:aims}

The aim of this project is to \textbf{design and develop an affordable, safe
and reliable ride-pooling platform for Dhaka city} that matches passengers
travelling in the same direction into shared rides, so that each passenger pays
a lower fare and vehicles are used more efficiently.

\section{Project Objectives}
\label{sec:objectives}

To achieve this aim, the project has the following objectives:

\begin{enumerate}
  \item To develop a \textbf{rider mobile app} where users can search for a
        ride, view available pools, join a pool, track their trip and pay.
  \item To develop a \textbf{driver mobile app} where drivers can go online,
        accept pools, navigate to pickup points and view their earnings.
  \item To build a \textbf{backend server} that handles authentication, ride
        requests, pool matching, fare calculation, payments and notifications.
  \item To implement a \textbf{fast location-based matching system} using the
        H3 hexagonal grid to find nearby pools with similar routes.
  \item To implement a \textbf{fair fare system} where the fare drops as more
        passengers share the ride (for example, 25\% off with two riders and
        35\% off with three).
  \item To support \textbf{local vehicles and their capacity} --- CAR
        (3 passengers) and CNG (2 passengers).
  \item To provide \textbf{safety features} such as female-only rides, rider
        ratings and the \textit{Priyo Sathi} (trusted companions) feature.
  \item To reduce misuse through a \textbf{cancellation penalty and cooldown
        system}.
  \item To support \textbf{advance (scheduled) booking} so riders can book
        trips ahead of time and be automatically grouped into a pool.
  \item To provide an \textbf{in-app wallet} and cash payment, with real-time
        notifications and in-app chat.
\end{enumerate}

\section{Project Challenges}
\label{sec:challenges}

During the development of RidePool we faced the following main challenges:

\begin{itemize}
  \item \textbf{Efficient matching:} Searching for compatible riders across a
        whole city had to be fast. We solved this using Uber's H3 hexagonal
        grid, which divides the map into small hexagon cells so nearby pools
        can be found with simple ring searches instead of heavy distance
        calculations.
  \item \textbf{Route similarity:} Two riders being close at pickup does not
        mean they are going the same way. We had to build a scoring system that
        combines route overlap, pickup distance, destination distance and how
        full the pool is.
  \item \textbf{Cost of map services:} Google Maps APIs are charged per
        request. We had to reduce the number of calls using caching, a single
        shared route request per search, and a simple geometric fallback when
        map calls fail or are skipped.
  \item \textbf{Race conditions:} When several riders join the same pool at
        once, the vehicle could become overfilled. We handled this by moving
        critical operations (joining a pool, accepting a pool, debiting a
        wallet) into atomic database functions with row locking.
  \item \textbf{Real-time updates:} Riders and drivers need to see changes (a
        new passenger, driver location, trip status) immediately. We used
        Supabase Realtime together with a polling fallback.
  \item \textbf{Scheduled rides:} Advance booking needed time windows, rider
        confirmation and automatic assignment, all without breaking the
        existing instant-ride flow.
  \item \textbf{Two mobile apps and one server:} Keeping data types consistent
        between the rider app, the driver app and the server was difficult, so
        we wrote the common data model (rides, pools, users, payments) into a
        separate \texttt{shared} types package as one reference for all three.
\end{itemize}

\section{Contribution}
\label{sec:contribution}

The main contribution of this project is a complete, working ride-pooling
system made for the local context of Dhaka. It includes two mobile
applications, a backend server and a database, all working together. The key
areas of contribution are described below.

\subsection{Accessibility Focus}

RidePool is designed to make comfortable travel accessible to more people:

\begin{itemize}
  \item \textbf{Affordable fares:} Sharing a ride reduces the cost per person,
        which makes car and CNG travel affordable for students and daily
        commuters.
  \item \textbf{Local vehicle support:} Along with cars, the system supports
        CNG auto-rickshaws, which are cheaper and widely used in Dhaka.
  \item \textbf{Flexible payments:} Riders can pay using the in-app wallet or
        cash, and the wallet is designed around local mobile banking methods
        such as bKash. All prices are shown in BDT.
  \item \textbf{Safe for everyone:} The female-only option and the Priyo Sathi
        feature make it easier for women and cautious riders to use shared
        rides.
  \item \textbf{Available on multiple platforms:} The rider app runs on
        Android, iOS and the web from a single codebase.
\end{itemize}

\subsection{Integration of Modern Tools}

The project combines several modern technologies:

\begin{itemize}
  \item \textbf{React Native with Expo} for building both mobile apps from one
        codebase, with Expo Router for navigation, Zustand for state management
        and NativeWind (Tailwind CSS) for styling.
  \item \textbf{Node.js, Express and TypeScript} for the backend server, with
        Zod for input validation.
  \item \textbf{Supabase (PostgreSQL)} for the database, user authentication
        and real-time updates.
  \item \textbf{H3 (by Uber)} for hexagon-based geospatial indexing and fast
        pool matching.
  \item \textbf{Google Maps APIs} for routes, distances and ETAs, with Leaflet
        maps for the web version.
  \item \textbf{Redis or in-memory caching} to improve speed and reduce map API
        costs.
  \item \textbf{Vitest and Jest} for automated testing of the server and apps.
\end{itemize}

\subsection{User-Centered Design}

The system is built around the needs of its users:

\begin{itemize}
  \item \textbf{Riders} can see a list of available pools with the fare, pickup
        distance, number of passengers and estimated time, and choose the one
        that suits them best.
  \item Riders can set a \textbf{gender preference}, save frequently used
        places, add trusted companions (Priyo Sathi), chat with co-riders and
        the driver, and rate each other after a trip.
  \item Riders who know their travel time in advance can \textbf{schedule a
        ride}, and the system will group them automatically without any
        searching.
  \item \textbf{Drivers} have a simple app to go online, accept pools, follow
        the pickup route and track their earnings.
  \item Users receive \textbf{notifications} at every important step, such as a
        pool being found, a driver arriving or a trip being completed.
  \item The \textbf{cancellation penalty} (a short cooldown after repeated
        deliberate cancellations) protects other riders in the pool from being
        affected by one person.
\end{itemize}

\subsection{Modular and Scalable Architecture}

RidePool follows a clean and modular structure so that it is easy to maintain
and extend:

\begin{itemize}
  \item The project is split into separate packages: \textbf{Server},
        \textbf{CarPoolApp} (rider), \textbf{DriverApp} and \textbf{shared}
        (common data types).
  \item The server follows a layered design: \textbf{Routes $\rightarrow$
        Controllers $\rightarrow$ Services $\rightarrow$ Database}. Each layer
        has one responsibility, so business logic (matching, fares, wallet,
        penalties) is kept separate from HTTP handling.
  \item Each feature (pools, rides, payments, wallet, ratings, messaging,
        advance booking) has its own route, controller and service files, so
        new features can be added without changing existing ones.
  \item Important settings such as search radius, time windows and cache mode
        are kept in a central configuration file and environment variables.
  \item The system can run in a lightweight \textbf{MVP mode} with in-memory
        caching, and can switch to \textbf{Redis} when more users need to be
        served.
  \item Security is handled through authentication middleware, rate limiting,
        input sanitization and database row-level security.
\end{itemize}

\section{Organization of the Report}
\label{sec:organization}

The rest of this report is organized as follows:

\begin{itemize}
  \item \textbf{Chapter 2 -- Literature Review:} Discusses existing
        ride-sharing and ride-pooling systems and related work, and compares
        them with RidePool.
  \item \textbf{Chapter 3 -- Project Description:} Describes the requirements,
        tools and resources, and the societal, environmental, ethical,
        feasibility, risk and economic aspects of the project.
  \item \textbf{Chapter 4 -- Project Analysis:} Explains the system
        architecture, how the app modules work together, the UML and database
        diagrams, the user interface, and the evaluation and testing of the
        system.
  \item \textbf{Chapter 5 -- Project Implementation:} Describes how the
        server, apps and database were built, how the technology stack is
        deployed, and how to set up and run the system.
  \item \textbf{Chapter 6 -- Result Analysis:} Presents the results of the
        working system and discusses how well it meets its aims.
  \item \textbf{Chapter 7 -- Conclusion:} Summarizes the work done, discusses
        the limitations of the project and suggests future improvements.
\end{itemize}
